\subsection{取值于 g 的连续多项式}

设 I 是一个有限集，令 \(\mathbf{P}(\mathfrak{g}^{\mathrm{I}};\mathfrak{g})\)（分别令 \(\hat{\mathbf{P}} (\mathfrak{g}^{\mathrm{I}};\mathfrak{g}))\)）为定义在 \(\mathbf{g}^{\mathrm{I}}\) 上且取值于 \(\mathfrak{g}\) 的连续多项式（分别为带连续分量的形式幂级数）的向量空间。回忆（《微分流形与解析流形》，R，附录），\(\mathbf{P}(\mathfrak{g}^{\mathrm{I}};\mathfrak{g})\) 带有 \(\mathbf{N}^{\mathrm{I}}\) 型分次，而 \(\hat{\mathbf{P}} (\mathfrak{g}^{\mathrm{I}};\mathfrak{g})\) 被认作向量空间 \(\mathbf{P} (\mathfrak{g}^{\mathrm{I}};\mathfrak{g})\) 的完备化，其中的拓扑由与 \(\mathbf{P}(\mathfrak{g}^{\mathrm{I}};\mathfrak{g})\) 的分次相关的滤子定义。此外，\(\mathbf{P}(\mathfrak{g}^{\mathrm{I}};\mathfrak{g})\) 是一个分次李代数，其括号对 \([f,g](\pmb {x}) = [f(\pmb {x}),g(\pmb {x})]\)、\(f, g\) 定义为 \(\mathrm{P}(\mathfrak{g}^{\mathrm{I}}; \mathfrak{g}), x \in \mathfrak{g}^{\mathrm{I}}\)；这个李代数结构可由连续性延拓到 \(\hat{\mathrm{P}}(\mathfrak{g}^{\mathrm{I}}; \mathfrak{g})\)，使其成为完备豪斯多夫滤李代数。

根据 § 6 第 3 号的命题 2，存在唯一的连续李代数同态 \(\phi_{\mathrm{I}}:u\mapsto \tilde{u}\)，它从 \(\hat{\mathbf{L}} (\mathbf{I})\) 到 \(\hat{\mathbf{P}} (\mathbf{g}^{\mathrm{I}};\mathfrak{g})\)，并将指标为 \(i\) 的不定元映到 \(\mathsf{pr}_{\mathrm{i}}\)（对所有 \(i\in \mathbb{I}\)），因为 \(\mathsf{pr}_{\mathrm{i}}\in \mathsf{P}(\mathfrak{g}^{\mathrm{I}};\mathfrak{g})\)。由此，当 \(\tilde{u}\in \mathsf{P}(\mathfrak{g}^{\mathrm{I}};\mathfrak{g})\) 时，\(u\in \mathbf{L}(\mathbf{I})\)；更确切地说，当 \(u\in \mathbf{L}(\mathbf{I})\) 时，\(\tilde{u}\) 正是 § 2 第 4 号中的多项式映射 \((t_i)\mapsto u((t_i))\)。另一方面，显然 \(\phi_{\mathrm{I}}\) 与 \(\mathbf{L}(\mathbf{I})\) 和 \(\mathbf{P}(\mathfrak{g}^{\mathrm{I}};\mathfrak{g})\) 的多重分次相容。若 \(u = \sum_{\nu \in \mathbf{N}^{\mathrm{I}}}u_{\nu}\)，其中对 \(u_{\nu}\in \mathbf{L}^{\nu}(\mathbf{I})\) 有 \(\nu \in \mathbf{N}^{\mathrm{I}}\)，则

\[
\tilde {u} = \sum_ {\nu \in \mathbf {N} ^ {\mathrm{I}}}  \tilde {u} _ {\nu}, \quad \text { where } \tilde {u} _ {\nu} \in \mathrm{P} _ {\nu} (\mathfrak {g} ^ {\mathrm{I}}; \mathfrak {g}).
\]

令 \(\pmb{u} = (u_j)_{j\in J}\) 为 \(\hat{\mathbf{L}} (\mathbf{I})\) 中元素的有限族，令 \(v\in \hat{\mathbf{L}} (\mathbf{J})\)，并令 \(w = v\circ \pmb {u}\)（§ 6，第 3 号）。记 \(\tilde{\pmb{u}} = (\tilde{u}_j)_{j\in \mathbb{J}}\in \mathfrak{J}\)。则

\[
\tilde {v} \circ \tilde {\boldsymbol {u}} = (v \circ \boldsymbol {u}) ^ {\sim}.\tag{2}
\]

这是通过连续延拓 § 6 第 3 号的公式 (7)，并应用《微分流形与解析流形》R 附录第 6 号得到的。

